\documentclass[12pt]{letter}

\usepackage[utf8]{inputenc}
\usepackage[T1]{fontenc}

\begin{document}

\begin{letter}{}

    \opening{Dear Editor,}
    
    Please find enclosed our manuscript entitled ``Bi-Level Degradation-Aware Planning of Shared Battery Energy Storage Systems for Coordinated TSO--DSO Operation'', which we wish to submit for consideration for publication in the \textit{Journal of Energy Storage}.
    
    This manuscript proposes a degradation-aware planning framework for shared battery energy storage systems installed at transmission--distribution interface nodes and jointly operated under coordinated TSO--DSO operation. The proposed approach is formulated as a bi-level stochastic optimization model in which the upper level determines siting, sizing, and staged investment decisions under investment-cost uncertainty, while the lower level evaluates these decisions through coordinated operation. To preserve tractability, the framework combines Benders' decomposition for long-term planning with an ADMM-based decentralized coordination mechanism for short-term operation.
    
    The main contribution of the paper is an integrated framework that explicitly links long-term shared ESS planning with coordinated transmission--distribution operation while accounting for battery degradation. In particular, the manuscript combines four elements that are rarely addressed simultaneously in the literature:
    \begin{itemize}
        \item siting and sizing of shared ESSs at TN--DN interface nodes;
        \item explicit representation of cycling-induced degradation;
        \item coordinated AC-feasible TSO--DSO operation; and
        \item a decentralized solution approach compatible with institutional separation and data privacy constraints.
    \end{itemize}
    
    The proposed framework is evaluated on integrated IEEE transmission--distribution test systems over a 15-year planning horizon. The results show that coordinated operation with shared ESSs can significantly reduce operating costs and renewable curtailment, while also eliminating voltage violations observed under uncoordinated operation. The analysis further shows that battery degradation materially affects ESS valuation and that planning-horizon discretization can influence siting and sizing decisions.
    
    We believe this manuscript is well aligned with the scope of the \textit{Journal of Energy Storage}, as it addresses the planning and operational value of battery energy storage systems in future power systems with high renewable penetration, while also highlighting the importance of degradation-aware modeling in long-term decision-making.
    
    All authors have approved the manuscript and agree with its submission to the \textit{Journal of Energy Storage}.
    
    Thank you for considering our work for publication in the \textit{Journal of Energy Storage}. We believe the manuscript will be of interest to researchers and practitioners working on energy storage planning and operation.
    
    \closing{Yours sincerely,}    
    Micael Simões\\
    On behalf of all authors\\
    INESC TEC / FEUP\\
    micael.f.simoes@inesctec.pt

\end{letter}

\end{document}